\chapter{Kafka: Mô hình lõi}
\label{ch:kafkacore}

Kafka là công nghệ đáng học sâu nhất trong stack này: nó xuất hiện trong
một phần lớn các mô tả công việc backend, và mô hình của nó thực sự khác
với một hàng đợi thông điệp (message queue). Chương này là Kafka độc lập
với nền tảng; Chương~11 là phía Spring.

\section{Một log, không phải một hàng đợi}

Trừu tượng lõi của Kafka là một \emph{log chỉ ghi nối thêm} (append-only
log). Một \emph{topic} (ví dụ \code{ts.user.registered}) được chia thành
các \emph{phân vùng} (partition); mỗi partition là một dãy \emph{record}
có thứ tự và bất biến, mỗi record ở một \emph{offset} được đánh số.
Record không bị xóa khi đọc --- chúng hết hạn theo chính sách retention
(dự án này: 168 giờ). Đọc không phá hủy; consumer chỉ nhớ \emph{mình đã
đọc tới đâu}.

\begin{figure}[h]
\centering
\begin{tikzpicture}[node distance=3mm and 8mm]
  \node[lbl] (t) at (0,1.35) {topic ts.user.registered};
  \foreach \p/\y in {0/0.6, 1/0} {
    \node[lbl,anchor=east] at (0.1,\y) {p\p};
    \foreach \i in {0,...,6} {
      \node[draw,minimum size=5mm,font=\scriptsize] at (0.65+\i*0.62,\y) {\i};
    }
  }
  \draw[flow] (5.2,1.15) node[lbl,right]{producer ghi nối} -- (4.7,0.72);
  \draw[flow] (2.55,-0.75) node[lbl,below]{offset của group: 4} -- (3.13,-0.28);
\end{tikzpicture}
\caption{Hai partition; producer ghi nối thêm ở đầu, consumer group nhớ
một offset cho mỗi partition.}
\end{figure}

Các hệ quả, mà một hàng đợi không có cái nào:

\begin{itemize}
  \item \textbf{Phát lại được.} Một consumer mới có thể bắt đầu từ offset
  0 và thấy lịch sử. Sang năm triển khai một service phân tích mới; nó
  có thể xử lý các lượt đăng ký của tuần trước.
  \item \textbf{Nhiều người đọc.} Notification-service và course-service
  cùng đọc sự kiện auth một cách độc lập, theo nhịp riêng, không giành
  của nhau --- mỗi \emph{consumer group} giữ offset riêng của mình.
  \item \textbf{Có thứ tự trong từng partition.} Các record có cùng
  \emph{key} băm về cùng một partition và giữ nguyên thứ tự. Đây là lý do
  các publisher trong dự án này dùng \code{userId} làm key: mọi sự kiện
  về một người dùng đến theo đúng trình tự. Giữa các partition thì không
  có thứ tự --- chỉ thiết kế dựa trên thứ tự theo key.
\end{itemize}

\section{Consumer group}

Một consumer group là một đội có tên chia nhau việc: mỗi partition được
gán cho đúng một consumer \emph{trong group}. Hai consumer trong group
\code{notification-group} trên một topic 2 partition nhận mỗi bên một
partition; consumer thứ ba ngồi không --- \emph{partition là đơn vị của
tính song song}. Khi một consumer tham gia hoặc chết, group
\emph{rebalance}: partition được gán lại và các thành viên tạm dừng
trong chốc lát. Các group khác nhau vô hình với nhau, và đó là cách cùng
một sự kiện tỏa ra cả \code{notification-group} lẫn \code{course-group}.

\section{Broker, replication, KRaft}

Một cụm production chạy nhiều \emph{broker}; mỗi partition có một leader
trên một broker và các \emph{replica} trên những broker khác
(\code{replication-factor}, RF). Producer ghi vào leader; \code{acks=all}
chờ cho tới khi các replica đồng bộ (ISR) đã có bản ghi. Việc điều phối
metadata trước đây cần ZooKeeper, một hệ phân tán thứ hai; Kafka hiện
đại thay nó bằng \emph{KRaft} --- các broker tự chạy một quorum Raft nội
bộ. Dự án này chạy KRaft với một node đóng cả hai vai \code{broker} và
\code{controller}, RF\,=\,1: ổn cho homelab, khả năng chịu lỗi bằng
không --- một đĩa hỏng là mất log.

\section{Mê cung advertised listener}

Lỗi Kafka được Google nhiều nhất. Từ \code{docker-compose.yml}:

\begin{lstlisting}[language=yaml]
KAFKA_LISTENERS: PLAINTEXT://:19092,CONTROLLER://:9093,EXTERNAL://:9092 (*@\co{1}@*)
KAFKA_ADVERTISED_LISTENERS: PLAINTEXT://kafka:19092,EXTERNAL://localhost:9092 (*@\co{2}@*)
KAFKA_CONTROLLER_QUORUM_VOTERS: 1@kafka:9093                (*@\co{3}@*)
\end{lstlisting}

\begin{description}
  \item[\co{1}] \emph{Listener}: các socket mà broker mở. Ba cái: một cho
  client bên trong mạng Docker, một cho controller KRaft, một công bố ra
  máy host.
  \item[\co{2}] \emph{Advertised} listener: các địa chỉ mà broker
  \emph{bảo client hãy dùng}. Đây là cái bẫy: yêu cầu đầu tiên của một
  Kafka client là ``leader của các partition ở đâu?'', và broker trả lời
  bằng những chuỗi này. Container phải được bảo là \code{kafka:19092}
  (phân giải được trên mạng Docker); công cụ trên laptop của bạn phải
  được bảo là \code{localhost:9092}. Một broker, hai đối tượng, hai địa
  chỉ advertised.
  \item[\co{3}] Quorum KRaft gồm một thành viên: node 1 là toàn bộ cuộc
  bầu cử.
\end{description}

\begin{pitfall}[kết nối được, rồi timeout]
Triệu chứng: client kết nối tới địa chỉ bootstrap thành công, rồi mọi
send/poll đều thất bại với \code{Connection to node -1 could not be
established} hoặc treo. Nguyên nhân: bootstrap chạy được, nhưng
\emph{câu trả lời metadata} advertise một địa chỉ mà client không với
tới (ví dụ laptop của bạn bị bảo là \code{kafka:19092}). Địa chỉ
bootstrap chỉ mở đầu cuộc trò chuyện; advertised listener mang toàn bộ
lưu lượng sau đó. Chẩn đoán bằng cách hỏi broker advertise cái gì ---
ví dụ \code{kcat -L -b localhost:9092} --- và kiểm tra đúng những chuỗi
đó có phân giải được từ nơi client chạy hay không.
\end{pitfall}

\section{Ngữ nghĩa chuyển giao}

Mặc định của Kafka là \emph{ít nhất một lần} (at-least-once): producer
gửi lại sau khi mất ack, hoặc consumer sập sau khi xử lý nhưng trước khi
commit offset, cả hai đều gây trùng lặp --- chứ không mất. Hãy thiết kế
consumer \emph{lũy đẳng} (idempotent: xử lý cùng một sự kiện hai lần là
vô hại): notification-service gửi trùng một email là chấp nhận được; một
service thanh toán thì sẽ phải theo dõi ID các sự kiện đã xử lý. Đúng
một lần (exactly-once) có tồn tại (producer lũy đẳng + transaction)
nhưng trả giá bằng độ phức tạp; hãy biết rằng ít nhất một lần + consumer
lũy đẳng là món cơm bữa của cả ngành.

\section{Ngoài phạm vi dự án}

Kafka được quản lý (AWS MSK, Confluent Cloud) bỏ đi việc vận hành broker
và bán cùng giao thức đó --- client của bạn đổi một chuỗi bootstrap và có
thêm SASL/TLS. Những cái tên khác trong mảng này giải quyết những bài
toán khác: RabbitMQ là một \emph{hàng đợi} định tuyến thông minh (phân
phối công việc, ack theo từng thông điệp, không phát lại); Redpanda
hiện thực lại giao thức Kafka mà không cần JVM. Khi Chương~25 chuyển
broker này vào cụm dưới Strimzi, không gì trong chương này thay đổi ---
Strimzi là \emph{vận hành} cho cùng một Kafka.

\begin{quiz}
\begin{enumerate}
  \item Các sự kiện đăng ký của một người dùng phải được xử lý theo thứ
  tự, nhưng khả năng mở rộng cũng quan trọng. Bạn dùng gì làm key, bao
  nhiêu partition, và giới hạn song song là gì?
  \item Hai group, \code{notification-group} và \code{course-group}, cùng
  đọc \code{ts.user.activated}. Giải thích cách cả hai nhận được mọi sự
  kiện.
  \item Phân biệt \code{LISTENERS} với \code{ADVERTISED\_LISTENERS}, và
  giải thích vì sao lỗi chúng gây ra xuất hiện \emph{sau} một kết nối
  thành công.
  \item Vì sao chuyển giao ít nhất một lần đẩy tính lũy đẳng về phía
  consumer, và chiến lược lũy đẳng hợp lý của notification-service sẽ là
  gì?
\end{enumerate}
\end{quiz}
